\section{Strengths and Weaknesses}

\subsection{Strengths}
\begin{enumerate}[\quad\bf1.]        %%列出1、2、3点
    \item \textbf {Data processing.} We process the data by the algorithm of min-max normalization to transfer all selling period of different length to 14-day-long one. It facilitates our model design and result calculation.
    \item \textbf {Visualization.} We vividly and diversely visualize the waybill distribution by violin plots and use line chart to show the short-term forecasts. Additionally, we integrate line chart and bar chart together to display the evolution trend of the sales volume and the freight rate to further investigate the relationship between them.
    \item \textbf {Comprehensiveness.} We consider enough factors during our model design and use Entropy Weight Method and AHP to determine the factor weight, which is relatively accurate.
\end{enumerate}
\subsection{Weaknesses}
While our model and strategy get higher revenue, there remain several types of model weaknesses:
\begin{enumerate}[\quad\bf1.]  %%列出1、2、3点
    \item \textbf {Subjectivity.} The calculation of some synthesized variables is subjective, such as using AHP to get the weight.
    \item \textbf {Lack of practice.} Our model is set under data of only ten flow directions. Since the model needs more data to be proved, it is possible for it to fail in some other flow directions.
\end{enumerate}